\documentclass[11pt]{article}

\usepackage[a4paper,margin=1in]{geometry}
\usepackage{amsmath,amssymb}
\usepackage{hyperref}
\usepackage[numbers]{natbib}
\usepackage{enumitem}

\title{M2 Internship Project\\
Discrete Diffusion Models and Transportation Flows}
\author{}
\date{}

\begin{document}
\maketitle

\section{Context and motivation}
Diffusion models and related methods such as flow matching are now central tools for generative modeling in continuous spaces (images, audio, video), where Euclidean geometry provides natural transportation mechanisms \citep{ho2020ddpm,lipman2022flowmatching}.
In contrast, for discrete data (notably text), autoregressive next-token models still dominate, partly because transportation notions are less canonical on finite spaces.
Designing meaningful costs, paths, and flows is harder, and existing discrete diffusion methods rely on specific noising kernels or expensive Markov simulation \citep{austin2021d3pm,campbell2022ctmc}.

This internship studies the mathematical interface between discrete diffusion and transport, with the objective of understanding learnability, sample complexity, and generalization, and of proposing principled extensions guided by optimal transport.

\section{Axis 1: Forward and inverse Markov processes}
The first axis analyzes forward noising chains and inverse Markov couplings used for generation.
In discrete DDPM-like models, the forward kernel is a key design choice (local kernels, masking/absorbing states, etc.) and directly controls reverse dynamics \citep{austin2021d3pm}.
Continuous-time jump-process formulations provide an even more flexible framework and tractable training objectives \citep{campbell2022ctmc}.

The plan is to begin with controlled binary settings on the Boolean hypercube, where denoising admits empirical-Bayes/score identities and Langevin-like sampling can be analyzed \citep{bach2025binary}, then extend to richer categorical and graph-structured spaces.

\section{Axis 2: Error propagation and generalization}
The second axis develops an end-to-end error decomposition for discrete diffusion pipelines.
The central question is how statistical estimation error, reverse-kernel approximation, and finite-step sampling errors combine in the final discrepancy between generated and target distributions.

In discrete spaces, the answer depends strongly on the discrepancy metric (e.g., total variation versus transport-based metrics). Recent work provides complementary tools: fine-grained error analyses for continuous diffusion \citep{hurault2025gaussian} and new guarantees for discrete diffusion samplers with improved vocabulary-size dependence \citep{liang2025discrete}.

The goal is to derive explicit or near-explicit rates in simple regimes, clarifying trade-offs between model complexity, data size, and sampler depth.

\section{Axis 3: Connections with optimal transport}
The third axis investigates how discrete diffusion relates to optimal transport.
In Euclidean spaces, diffusion is linked to entropy gradient flows via JKO \citep{jordan1998jko}, while modern flow methods can sometimes be interpreted through transport velocities, with known limitations \citep{hertrich2025rfot}.

On finite spaces, a rigorous transport geometry exists for Markov chains, where Markov semigroups are entropy gradient flows under adapted transport metrics \citep{maas2011gradient,mielke2011markov}.
The internship will leverage this framework to understand what transportation structures are implemented by practical discrete diffusion samplers and whether transport-informed designs can improve efficiency or theory.

\section{Expected outcomes}
\begin{itemize}[leftmargin=*]
\item A mathematically clear benchmark of forward/reverse kernel choices in discrete diffusion.
\item A decomposition of generation error into estimation, approximation, and sampling terms.
\item Concrete links (and limits) between discrete diffusion and transport geometry.
\item Practical guidance for noising-kernel and reverse-coupling design.
\end{itemize}

\bibliographystyle{plainnat}
\bibliography{references}

\end{document}
